\section{Conclusion}
\label{sec:conclusion}

This paper asked one question: can tensor cores be applied to ALS matrix factorization? Tensor cores sit inside modern GPUs, including ordinary gaming cards, and they do dense matrix multiplication very fast, but a recommendation system deals with a sparse matrix, and a sparse matrix does not feed tensor cores well. The answer built here rests on sorting. Users are sorted so that the ones who give the most ratings come first, and items so that the most popular come first. After the sorting, the dense part of the matrix gathers in one corner, and on that corner tensor cores become usable. APR-BALS wraps this idea in a complete training system: a three-tier dispatch sends the dense work to the tensor cores in FP16 with FP32 accumulation, the sparse tail goes to a register-tiled scalar FP32 kernel, and a blocked batched Cholesky solver keeps the $K^{3}$ solve stage from eating the gain. The baseline and the accelerated path share one binary, so the comparison stays fair. On an RTX~3060 the Netflix Prize dataset trains $4.1\times$ to $14.8\times$ faster than the FP32 baseline over $K = 16$--$96$ under the identical protocol, the test RMSE agrees to within run-to-run noise ($\le 0.005\%$), and the results repeat on MovieLens 10M, 20M, and 32M and, under the identical protocol, on a Tesla T4 ($10.8\times$ on Netflix).

The method has clear borders. The supported latent dimensions are $K = 16, 32, 48, 64, 96$ only, enforced by a compile-time guard; adding another value needs a new WMMA template instantiation plus a register-spill check. At high $K$, the tiled Cholesky solve remains the largest single component of the accelerated iteration, $39\%$ of GPU time in the fully optimized $K=96$ recipe of Section~\ref{sec:extended}, and phase-isolation probes indicate that kernel now sits near a latency-bound structural floor. The system optimizes explicit ratings and evaluates by RMSE; implicit-feedback weighting and ranking metrics are outside the current scope. Only a single GPU is used. The sparse path uses atomic additions, so an APR run repeats only to about $10^{-6}$; bitwise-reproducible training would require a deterministic reduction. Under the per-entity weighted regularizer that serves as the protocol of record, the held-out accuracy improves monotonically with $K$, so $K=96$ is simultaneously the fastest-accelerated and the most accurate configuration; Section~\ref{sec:extended} shows that under a plain uniform penalty the opposite holds, which is why the weighted term is the default and not an option.

Section~\ref{sec:extended} closes several questions this paper opened. Switching to weighted, per-entity regularization removes the high-$K$ overfitting turn entirely and, because it also caps the worst-case condition number of every entity's system, unlocks a narrower-precision Cholesky solver at $K=96$ that reproduces the FP32 solver's test RMSE exactly while cutting the solve stage by $23\%$; fusing the regularizer into the solver's existing load pass and relocalizing the periodic RMSE check onto the tensor-core path's own FP16 factor images remove the two next-largest overheads, a further $18\%$ off the $K=96$ wall time. Against cuMF\_als \cite{tan2016faster}, the most cited GPU implementation of conjugate-gradient ALS, APR-BALS wins per-iteration speed by $3.0\times$--$6.4\times$ on the same frozen split while agreeing with cuMF's train RMSE to within $0.33\%$. The comparison was then extended to the two other solver families BALS itself was evaluated against: against cuMF\_CCD, a GPU coordinate-descent solver \cite{yu2012scalable}, APR-BALS is $3.4\times$--$4.1\times$ faster, matching or exceeding the margin BALS's own paper reported over the identical baseline on a TITAN RTX; against cuMF\_SGD \cite{xie2017cumf}, which supports only rank $128$, APR-BALS at rank $96$ matches or beats its held-out accuracy at comparable wall time on every dataset large enough to be conclusive, despite using fewer factors. A faithful, independently reimplemented scalar-FP32 BALS baseline, complete with BALS's own frequency-based reordering and symmetric register-tile Gram mapping, in turn runs $1.6\times$--$2.1\times$ \textit{slower} than cuMF\_als on this Ampere GPU, the reverse of BALS's own pre-Ampere result and a reminder that a specialized kernel's advantage is tied to the hardware generation it was tuned against; APR-BALS's $10\times$-plus margin over that same faithful baseline isolates the tensor-core precision dispatch, rather than the tiled storage format it is built on, as the source of the speedup reported throughout this paper. Finally, extrapolating the ALS factor sequence with heavy-ball momentum, orthogonally to every kernel optimization above, cuts the sweep count to convergence from 20--25 to 15 at every rank on every dataset tested, worth a further $1.3\times$--$1.7\times$ wall time at equal or slightly better test accuracy; compounded, the shipped recipe trains Netflix at $K=96$ $26.0\times$ faster end-to-end than the equally optimized FP32 baseline. Two solver alternatives were also implemented and measured to close off the question of how far the Cholesky stage can be pushed: a vendor batched factorization (cuSOLVER) and a tensor-core batched-multiply formulation both ran slower, because at the ranks this system targets the per-entity matrices sit below the size at which tensor cores factor efficiently, so the shared-memory scalar Cholesky is retained by measurement, not by default. What remains open is a lean per-entity batched-GEMM Gram kernel that might let a pure-FP32, non-tensor-core path close the gap to cuMF on Ampere, though this would largely re-derive cuMF's own design; a residual-gated Woodbury solve for the lightest entities measured a real but modest further $4\%$ off the $K=96$ wall time and is parked as an algorithmic side note rather than folded into the headline recipe; and writing the Gram matrices directly in the solver's tiled layout would turn the load phase into a straight vectorized copy, a change the already-implemented dual-stream pipeline would convert into wall-clock gains automatically. BF16 and TF32 operand formats run on the same tensor cores, widen the exponent range, and would likely shrink the already small high-$K$ deltas further; extending the dispatch to the confidence-weighted implicit-feedback objective \cite{hu2008collaborative} would broaden applicability; and RMSE could be complemented with ranking metrics.

One last point goes beyond recommendation. Many sparse workloads hide dense per-entity substructure: ridge-regression sweeps, Gaussian-process approximations, graph feature aggregation. The same tiered-precision pattern has a fair chance of working there too, and the honest-baseline habit of this work, sharing one binary between the reference and the accelerated path, and re-deriving even a competitor's own baseline faithfully enough to see where its advantage does and does not survive a change of hardware, travels even further. Any acceleration claim, in any field, becomes stronger when it is measured this way.
